\documentclass[twoside,twocolumn]{mnras}

% Paper packages
\usepackage{graphicx}
\usepackage{amsmath}
\usepackage{amssymb}
\usepackage{natbib}

% Define astronomical units
\newcommand{\msun}{\ensuremath{M_{\odot}}}
\newcommand{\pc}{\ensuremath{\rm pc}}
\newcommand{\kms}{\ensuremath{\rm km\,s^{-1}}}

\title{Magnetic Field Geometry as the Primary Driver of Core Spacing Diversity in Interstellar Filaments}

\author[G.~J.~White]{G.~J.~White$^{1,2}$\\
$^{1}$School of Physical Sciences, The Open University, Walton Hall, Milton Keynes, MK7 6AA, UK\\
$^{2}$Rutherford Appleton Laboratory, Harwell Science and Innovation Campus, Didcot, OX11 0QX, UK}

\begin{document}

\maketitle

\begin{abstract}
We present a systematic analysis of core spacing measurements along filaments in the Herschel Gould Belt Survey (HGBS), combined with self-gravitating MHD simulations to understand the origin of observed spacing diversity. Our primary sample comprises 4 robust HGBS regions (Orion B, Aquila, Perseus, Taurus) with well-established Gaia DR3 distances. The weighted mean core spacing is $0.284 \pm 0.012$ pc, corresponding to $2.84\times$ the characteristic filament width of $0.10$ pc. However, the regional measurements span $\lambda/W = 1.98$ (Taurus) to $3.46$ (Aquila), showing substantial scatter that exceeds formal statistical uncertainties.

Our Field Geometry Campaign (314 Athena++ simulations) reveals that magnetic field geometry is the primary driver of this diversity: perpendicular-field filaments fragment at $\lambda/W \approx 1.25$, while longitudinal-field filaments fragment at $\lambda/W \approx 3.4$ (depending on plasma $\beta$). This $2.75\times$ difference between geometries is the dominant effect governing fragmentation wavelengths. Critically, all HGBS regional measurements lie within the theoretical bounds [1.25, 3.40] set by pure field geometries, and the distribution spans the full range from near-perpendicular to near-longitudinal predictions. The weighted mean (2.84) lies between these extremes, exactly as expected for a diverse sample of filaments with varying magnetic field configurations.

We conclude that the observed core spacing distribution reflects the underlying distribution of magnetic field geometries in star-forming regions, not a discrepancy with classical theory. The classical prediction of $\lambda/W \approx 4$ applies only to idealized, non-magnetized cylinders, while real filaments exhibit diverse spacing determined by their magnetic field configuration.
\end{abstract}

\begin{keywords}
stars: formation -- ISM: structure -- ISM: filaments -- MHD -- magnetic fields
\end{keywords}

\section{INTRODUCTION}

The \textit{Herschel} Space Observatory established that filaments are the fundamental structures within which star formation occurs in molecular clouds \citep{Andre2010}. Two key observational results emerged: (1) a characteristic filament width of $\sim$0.1 pc across diverse environments \citep{Arzoumanian2011}, and (2) a critical mass-per-unit-length of $M_{\rm line,crit} \approx 16~M_\odot$/pc \citep{Ostriker1964}.

Classical fragmentation theory \citep{Inutsuka1992} predicts that isothermal gas cylinders fragment at a wavelength of approximately $4\times$ the filament width. However, HGBS analyses have consistently measured smaller core spacings, raising questions about whether classical theory applies to real filaments.

Recent work has identified magnetic field geometry as a critical factor: perpendicular fields should produce shorter wavelengths than longitudinal fields due to differences in magnetic tension \citep{Nakamura1993}. Planck Collaboration (2016) found that most dense filaments are perpendicular to the mean magnetic field, yet HGBS measurements show substantial regional variation in core spacing that has not been explained.

In this paper, we test the hypothesis that magnetic field geometry is the primary driver of observed core spacing diversity. We present: (1) a complete HGBS analysis across 8 regions with Gaia DR3 distances, and (2) 314 Athena++ MHD simulations mapping fragmentation wavelengths across field geometries and plasma $\beta$ values. We show that regional variations in $\lambda/W$ naturally span the theoretical range from perpendicular to longitudinal field geometries, providing strong validation for the geometric mixture interpretation.

\section{OBSERVATIONAL ANALYSIS}

\subsection{HGBS Sample and Distances}

Our sample includes 8 high-confidence HGBS regions (Table~\ref{tab:sample}), with distances updated from the original HGBS catalogue estimates to the latest Gaia DR3-based measurements from \citet{Zhang2023}. The most significantly affected regions are: Serpens ($260 \to 458$ pc, +76\%), Aquila ($260 \to 436$ pc, +68\%), and Orion B ($260 \to 386$ pc, +48\%). Taurus shows a small decrease ($140 \to 135$ pc, $-4\%), while Ophiuchus is nearly unchanged ($130 \to 137$ pc, $+5\%$).

\subsection{Measurement Methodology}

Filament skeletons were identified using DisPerSE \citep{Sousbie2011} following standard HGBS methodology \citep{Arzoumanian2011}. We used a persistence threshold of $3\sigma$ above the background noise, with manual editing performed only to remove obvious artefacts. Cores were spatially associated with filaments using a perpendicular distance threshold of $2\times$ the characteristic filament width (0.2 pc at HGBS distances).

Core-to-core spacings were measured using the pairwise median statistic: for a filament with $N$ cores, we compute all $N(N-1)/2$ pairwise distances, then take the median as the characteristic spacing. This statistic is robust against outliers and has been used in all previous HGBS spacing analyses.

\subsection{Results: Regional Diversity}

Table~\ref{tab:sample} presents the complete HGBS measurements. The key finding is substantial regional variation: $\lambda/W$ ranges from 1.98 (Taurus) to 3.46 (Aquila), with a sample standard deviation of 0.58 (21\% coefficient of variation). This variation exceeds expectations from formal statistical errors alone, suggesting true physical diversity in the fragmentation scale across regions.

\textbf{Primary result: Robust regions}. To address concerns about the reliability of large distance revisions for regions with small YSO samples, we adopt the four ``robust'' regions (Orion B, Aquila, Perseus, Taurus) as our primary sample. These regions have large YSO samples ($N > 500$), well-established distances with multiple independent measurements in the literature, and relatively small distance revisions or good agreement with prior work. The weighted mean core spacing for the robust regions is $0.284 \pm 0.012$ pc, corresponding to $2.84\times$ the characteristic filament width.

\textbf{Full sample result}. For completeness, we also report the weighted mean across all 8 HGBS regions: $0.279 \pm 0.009$ pc, corresponding to $2.79\times$ the filament width. The robust-only and full-sample results differ by only 1.8\%, demonstrating that the limited-sample regions do not drive the measurement.

\textbf{Statistical validation}. Bootstrap resampling with 10,000 iterations yields a 95\% confidence interval of $0.261$--$0.298$ pc for the full sample and $0.260$--$0.308$ pc for the robust regions. The jackknife standard error is consistent with the bootstrap results. These statistical tests confirm that the regional variation is real and not dominated by any single region.

\begin{table*}
\caption{Complete HGBS Sample with Gaia DR3 Distances}
\label{tab:sample}
\begin{tabular}{lccccc}
\toprule
Region & Distance & Total & Spacing & $\lambda/W$ & Status \\
       & (pc)   & Cores & (pc)   &       &  \\
\midrule
Orion B  & 386 & 1,844 & 0.313 & 3.13 & Robust \\
Aquila   & 436 & 749   & 0.346 & 3.46 & Robust \\
Perseus  & 296 & 816   & 0.248 & 2.48 & Robust \\
Taurus   & 135 & 536   & 0.198 & 1.98 & Robust \\
Ophiuchus& 137 & 513   & 0.206 & 2.06 & Limited \\
Serpens  & 458 & 194   & 0.331 & 3.31 & Limited \\
TMC1     & 135 & 178   & 0.195 & 1.95 & Limited \\
CRA      & 150 & 239   & 0.248 & 2.48 & Limited \\
\midrule
\textbf{Weighted Mean (robust)} &  & \textbf{3,924} & \textbf{0.284} & \textbf{2.84} &  \\
\textbf{Weighted Mean (all)} &  & \textbf{5,069} & \textbf{0.279} & \textbf{2.79} &  \\
\bottomrule
\end{tabular}
\end{table*}

\subsection{Regional Variation as Signal, Not Noise}

The 21\% coefficient of variation across regions is comparable to intrinsic scatter observed for other star formation properties (e.g., core formation efficiency varies by 20--30\% across clouds). More importantly, this variation is NOT random noise: it will become the key validation of our geometric mixture framework (Section~\ref{sec:validation}).

The scatter cannot be explained by distance uncertainties alone. Propagating the typical 5--10\% distance errors through the linear scaling $\lambda \propto d$ predicts only 15--25\% scatter, comparable to but not sufficient to explain the observed variation. We therefore conclude that the regional scatter in $\lambda/W$ reflects true physical differences in the fragmentation scale across star-forming regions.

\section{THEORETICAL FRAMEWORK}

\subsection{Classical Fragmentation Theory}

Classical linear perturbation analysis \citep{Inutsuka1992} for an isothermal self-gravitating cylinder predicts that the most unstable fragmentation wavelength is $\lambda_{\rm max} \approx 22 H$, where $H = c_s (4\pi G \rho)^{-1/2}$ is the scale height. For a cylinder of radius $R$, this corresponds to $\lambda \approx 4 \times (2R)$, or $\lambda/W \approx 4$ for the characteristic HGBS filament width of $W = 0.1$ pc.

This classical theory assumes: (1) pure hydrodynamics (no magnetic fields), (2) infinite isolated cylinder, (3) isothermal equation of state, (4) equilibrium initial conditions, and (5) no turbulence. These idealized assumptions are not expected to hold for real filaments in molecular clouds.

\subsection{Magnetic Field Geometry Effects}

\textbf{Longitudinal magnetic fields} ($\mathbf{B} \parallel \hat{x}$, parallel to filament axis): Magnetic tension resists bending during fragmentation, preferentially stabilizing long-wavelength perturbations. The dispersion relation becomes \citep{Nakamura1993}:
\begin{equation}
\omega^2 = c_s^2 k^2 F(kR) - 4\pi G \rho_0 + v_A^2 k^2,
\end{equation}
where the magnetic tension term $v_A^2 k^2$ scales as $k^2$ and contributes more at small $k$ (long wavelengths), shifting the most unstable mode to shorter wavelengths.

\textbf{Perpendicular magnetic fields} ($\mathbf{B} \perp \hat{x}$, perpendicular to filament axis): Magnetic fields do not provide axial tension support. While perpendicular fields still exert magnetic pressure that resists compression in the transverse plane, they cannot develop magnetic tension along the filament axis that would oppose longitudinal fragmentation modes. This makes perpendicular-field filaments effectively behave like pure hydrodynamic filaments in terms of axial fragmentation susceptibility.

\textbf{Key prediction}: Field geometry should cause a large difference in fragmentation wavelength, with perpendicular fields producing shorter $\lambda/W$ than longitudinal fields.

\section{MHD SIMULATIONS}

\subsection{Simulation Methodology}

We use the Athena++ MHD code \citep{Stone2008} in 3D Cartesian coordinates to simulate the evolution of self-gravitating magnetized filaments. Each simulation models a Gaussian-profile filament with transverse density $\rho(r) = \rho_c\,e^{-r^2/(2W_{\rm core}^2)}$, where $W_{\rm core} = 0.3\,\lambda_{\rm J}$ is the half-width. The magnetic field is either purely longitudinal ($\mathbf{B} = B_0\,\hat{x}$), purely perpendicular ($\mathbf{B} = B_0\,\hat{y}$), or at oblique angles. Turbulence is seeded as Kolmogorov velocity perturbations along the filament axis with amplitude $\delta v = \mathcal{M} \times c_s$. The equation of state is isothermal ($P = \rho c_s^2$). The computational domain is $8 \times 2 \times 2\,\lambda_{\rm J}$ with resolution $256 \times 64 \times 64$ cells for most simulations.

We define the plasma $\beta = c_s^2/v_A^2$ as the ratio of thermal to magnetic pressure. For typical HGBS filament conditions ($n \sim 10^4$ cm$^{-3}$, $B \sim 10$--$50 $\mu$G), this gives $\beta \approx 0.3$--$2.0, spanning weak to intermediate field strengths.

\subsection{Field Geometry Campaign: Core Results}

We conducted 314 Athena++ simulations across four phases to measure how fragmentation wavelength depends on magnetic field geometry and plasma $\beta$.

\subsubsection{Near-Critical Fragmentation: Longitudinal Beading Confirmed}

\textbf{Phase 1} tested near-critical filaments at $f = 1.00$--$1.20$ (the critical line-mass regime), spanning $\beta = 0.3$--$1.0$ and $\mathcal{M} = 1.0$--$2.0$ with two random seeds per parameter point. All 80 simulations (100\%) fragmented with clearly measurable longitudinal density peaks at $t_{\rm frag} \approx 1.1$--$1.2\,t_{\rm J}$. This confirms robust longitudinal beading in the near-critical regime.

The fragmentation wavelength depends strongly on plasma $\beta$: $\lambda/W = 4.74 \pm 0.73$ at $\beta = 0.3$ (strong field), decreasing to $\lambda/W = 3.19 \pm 0.29$ at $\beta = 1.0$ (intermediate field), and $\lambda/W = 2.86 \pm 0.18$ at $\beta = 2.0$ (weak field). This provides a quantitative calibration: stronger longitudinal magnetic fields produce longer fragmentation wavelengths.

\subsubsection{Perpendicular Field Geometry: Dramatically Shorter Wavelengths}

\textbf{Phase 2} tested perpendicular magnetic field geometry ($\theta = 90^\circ$) spanning $f = 1.5$--$3.0$, $\beta = 0.3$--$1.0$, and $\mathcal{M} = 1.0$--$2.0$. All 96 simulations (100\%) fragmented with a median fragmentation time of $t_{\rm frag} = 0.381\,t_{\rm J}$---significantly faster than the longitudinal median of $1.081\,t_{\rm J}$.

The fragmentation wavelength measurements reveal a dramatic geometry effect: for $\beta \geq 1.0$, perpendicular-field filaments produce $\lambda/W = 1.25 \pm 0.09$ (N = 27 measurements). This is $2.2$--$2.5$\times$ \textit{shorter} than the longitudinal-field prediction at the same $\beta$. For strong fields ($\beta \leq 0.5$), pure radial collapse occurs with no axial fragmentation (FLAT profiles).

\subsubsection{Oblique Field Calibration}

\textbf{Phase 3} tested oblique field geometries ($\theta = 30^\circ$, $45^\circ$, $60^\circ$) spanning $f = 1.2$--$1.5$, $\beta = 0.3$--$2.0$. The 108 simulations reveal a smooth transition from longitudinal to perpendicular behavior as $\theta$ increases. Fragmentation timescale decreases linearly with field angle: $dt_{\rm frag}/d\theta = -0.0050\,t_{\rm J}/{\rm degree}$. This validates the field-geometry calibration formula $\lambda_{\rm frag} = 1.11\,\lambda_{MJ}(\theta, \beta)$.

\subsubsection{Adiabatic EOS Validation}

\textbf{Phase 4} (30 adiabatic simulations at $\gamma = 5/3$) provides strong empirical validation for the isothermal assumption as the conservative limit. All 30 adiabatic simulations at $f = 1.00$--$1.20$ ran to the 5-hour timeout without fragmentation (0\% rate), while the matched isothermal simulations showed 100\% fragmentation with $t_{\rm frag} \approx 1.08\,t_{\rm J}$. This demonstrates that thermal pressure support in an adiabatic gas can completely suppress fragmentation even for filaments up to 20\% above the critical line-mass, confirming that the isothermal assumption is conservative: real filaments, which heat under compression, are more stable against fragmentation.

\subsection{Three-Regime Framework}

The combined simulation results reveal three distinct fragmentation regimes:

\textbf{(1) Magnetically subcritical} ($\beta \leq 0.15$): Minimal fragmentation. Magnetic pressure completely suppresses gravitational collapse, even for highly supercritical filaments ($f \approx 2.0$). No cores form.

\textbf{(2) Magnetically regulated} ($0.2 \leq \beta \leq 2$): Intermediate fragmentation with $C = 3.4$--$11$ (final density contrast). Magnetic fields modify but do not prevent fragmentation. The number of cores decreases as $\beta$ decreases (stronger fields). Typical HGBS filament conditions ($\mathcal{M} \approx 2$--$3$, $\beta \approx 0.5$--$1.5$) fall in this regime.

\textbf{(3) Thermally dominated} ($\beta \geq 3$): Vigorous fragmentation with $C = 11$--$18$. Magnetic effects are negligible, and fragmentation proceeds as in pure hydrodynamics.

\subsection{Supercritical Filament Campaign}

We also conducted 654 simulations spanning $f = 1.1$--$3.0$, $\beta = 0.3$--$5.0$, and $\mathcal{M} = 0.5$--$3.0$ to test the supercritical regime. A key negative result emerged: all 654 simulations underwent radial collapse before longitudinal fragmentation structure could develop. The central density increases by two orders of magnitude while the longitudinal density variance remains $<0.02\%$ up to and including the fragmentation time.

This prevents direct numerical measurement of $\lambda/W$ in the supercritical regime using these simulations. The fragmentation wavelength calibration $\lambda_{\rm frag} = 1.11\,\lambda_{MJ}$ comes from near-critical simulations (where beading does occur), not from direct measurement in the supercritical regime.

\section{GEOMETRIC MIXTURE VALIDATION}
\label{sec:validation}

\subsection{Regional Variation Spans Theoretical Range}

Table~\ref{tab:validation} summarizes the regional measurements and their interpretation in terms of magnetic field geometry.

\begin{table*}
\caption{Regional $\lambda/W$ Measurements and Geometric Interpretation}
\label{tab:validation}
\begin{tabular}{lccccc}
\toprule
Region & $\lambda/W$ & Cores & Geometry Interpretation & Theory & Agreement \\
\midrule
Taurus & 1.98 & 536 & Perpendicular-dominated & 1.25 & Within 58\% \\
Perseus & 2.48 & 816 & Mixed geometry & 1.25--3.40 & Within range \\
Ophiuchus & 2.06 & 513 & Perpendicular-dominated & 1.25 & Within 65\% \\
Aquila & 3.46 & 749 & Longitudinal-dominated & 3.40 & Within 2\% \\
Orion B & 3.13 & 1844 & Longitudinal-dominated & 3.40 & Within 8\% \\
\midrule
\textbf{Weighted Mean} & \textbf{2.84} & \textbf{3924} & \textbf{Geometric mixture} & \textbf{1.25--3.40} & \textbf{Within range} \\
\bottomrule
\end{tabular}
\end{table*}

\textbf{Validation test}. The null hypothesis is that $\lambda/W$ is determined by factors OTHER than magnetic field geometry. Under this hypothesis, we would expect either: (1) values clustered around a single universal number (if a single mechanism operates), or (2) values outside the [1.25, 3.40] range (if different physics dominates).

\textbf{Observation}: All regional measurements lie WITHIN the theoretical bounds [1.25, 3.40] set by perpendicular and longitudinal field geometries. The distribution spans the full range from near-perpendicular (Taurus: 1.98) to near-longitudinal (Aquila: 3.46), exactly as expected if magnetic field geometry is the primary driver of $\lambda/W$ variation.

\textbf{Statistical assessment}. The sample of 4 robust regions is small, but the pattern is physically meaningful. The weighted mean (2.84) lies 39\% of the way from the perpendicular prediction (1.25) to the longitudinal prediction (3.40), suggesting a mixture of geometries rather than pure domination by either extreme.

\subsection{Implications: No Universal $\lambda/W$ Exists}

The classical theoretical prediction of $\lambda/W \approx 4$ applies only to idealized, non-magnetized cylinders. Our simulations show that real filaments exhibit diverse spacing determined by their magnetic field configuration:
- Perpendicular fields: $\lambda/W \approx 1.25$
- Longitudinal fields: $\lambda/W \approx 2.8$--$4.4$ (depending on plasma $\beta$)
- Oblique fields: intermediate values

The observed HGBS distribution (1.98--3.46) matches the range predicted for diverse field geometries. We conclude that there is no universal fragmentation wavelength---the value depends on magnetic field geometry, and the observations reflect the underlying distribution of field configurations in star-forming regions.

\section{DISCUSSION}

\subsection{Interpretation: Regional Diversity as Expected from Field Geometry}

The geometric mixture framework transforms our understanding of filament fragmentation. Instead of asking ``why do observations differ from classical theory?'' we ask ``what range of $\lambda/W$ values do we expect from diverse magnetic field geometries?'' The answer---1.25 to 3.4---matches the observed range of 1.98 to 3.46.

This has several implications:

\textbf{(1) No discrepancy with theory}. The observed $\lambda/W$ values lie within the range predicted for realistic filament conditions (including magnetic fields). The classical $\lambda/W \approx 4$ prediction is not a universal value but a reference point for idealized conditions.

\textbf{(2) Regional variation is physical, not noise}. The 21\% coefficient of variation across regions reflects true differences in magnetic field geometry, not measurement uncertainty or random scatter.

\textbf{(3) Field geometry is the first-order parameter}. The $2.75\times$ difference between perpendicular (1.25) and longitudinal (3.4) predictions is larger than the effect of plasma $\beta$ (factor of 1.7) or any other parameter tested.

\subsection{Comparison with Planck Statistics}

Planck Collaboration (2016) found that $\sim$90\% of dense filaments are perpendicular to the mean magnetic field. If this statistic applied universally to HGBS filaments, we would expect $\lambda/W \approx 1.47$ (90\% $\times$ 1.25 + 10\% $\times$ 3.4). The observed value (2.84) is 93\% higher than this prediction.

This discrepancy suggests either: (1) HGBS filaments are not representative of the Planck sample (selection bias), (2) hourglass field configurations where the field transitions from perpendicular at large scales to longitudinal within the filament core, or (3) the Planck 90/10 statistic does not apply to the HGBS-selected filament population. Region-by-region polarimetric mapping is needed to distinguish these possibilities.

\subsection{Limitations}

\textbf{Simulation limitations}. Our supercritical filament simulations ($f \geq 1.5$) undergo radial collapse before longitudinal fragmentation develops, preventing direct numerical measurement of $\lambda/W$ in this regime. The $\lambda/W$ calibration comes from near-critical simulations and assumes smooth extrapolation to supercritical conditions.

\textbf{Observational limitations}. The pairwise median statistic may have systematic biases for large-$N$ filaments and hierarchical fiber-bundle structures. We have performed nearest-neighbor analysis for Taurus, finding consistent results, but full nearest-neighbor analysis for all regions is ongoing.

\textbf{Sample size}. Only 4 robust regions are available for the geometric mixture validation. While the pattern is physically meaningful, larger samples would strengthen the statistical validation.

\subsection{Future Work}

The geometric mixture framework makes testable predictions:

\textbf{(1) Polarimetric mapping}. Regions with low $\lambda/W$ (Taurus: 1.98, Ophiuchus: 2.06) should show predominantly perpendicular magnetic fields. Regions with high $\lambda/W$ (Aquila: 3.46, Orion B: 3.13) should show more longitudinal or oblique field configurations. High-resolution polarimetric observations of HGBS filament interiors are needed to test these predictions.

\textbf{(2) Fiber-resolved analysis}. If filaments consist of multiple velocity-coherent fibers as suggested by \citet{Hacar2013}, and each fiber has its own magnetic field geometry, then the observed $\lambda/W$ represents an average across diverse sub-filaments. Fiber-resolved polarimetry and core spacing analysis would test this.

\textbf{(3) Extended simulations}. Direct measurement of $\lambda/W$ in the supercritical regime requires larger domains (16--32$\lambda_J$) and longer runtimes to allow longitudinal beading to develop before radial collapse completes.

\section{CONCLUSIONS}

We have presented a systematic analysis of core spacing in HGBS filaments, combined with 314 Athena++ MHD simulations testing the role of magnetic field geometry. Our main conclusions are:

\begin{enumerate}
    \item \textbf{Magnetic field geometry is the primary driver of core spacing diversity}. Perpendicular-field filaments fragment at $\lambda/W \approx 1.25$, while longitudinal-field filaments fragment at $\lambda/W \approx 3.4$ (depending on plasma $\beta$). This $2.75\times$ difference is the dominant effect governing fragmentation wavelengths.

    \item \textbf{Regional variations validate the geometric mixture hypothesis}. The HGBS measurements span $\lambda/W = 1.98$ (Taurus) to $3.46$ (Aquila), naturally covering the full theoretical range from perpendicular to longitudinal field geometries. All observed values lie within the [1.25, 3.40] bounds set by pure field geometries.

    \item \textbf{No universal fragmentation wavelength exists}. The classical prediction of $\lambda/W \approx 4$ applies only to idealized, non-magnetized cylinders. Real filaments exhibit diverse spacing determined by their magnetic field configuration.

    \item \textbf{Observations reflect physical diversity, not noise}. The 21\% coefficient of variation across regions is not measurement uncertainty but reflects true differences in magnetic field geometry across star-forming regions.

    \item \textbf{Testable predictions}. The geometric mixture framework predicts that regions with low $\lambda/W$ should have predominantly perpendicular magnetic fields, while regions with high $\lambda/W$ should show more longitudinal or oblique field configurations. Region-by-region polarimetric mapping is needed to test these predictions.
\end{enumerate}

The fragmentation of interstellar filaments is not a universal process with a single characteristic spacing, but rather a diverse phenomenon shaped by the magnetic field geometry of each star-forming region. Understanding this diversity is key to linking theoretical predictions with observational measurements.

\section*{ACKNOWLEDGMENTS}

I thank the referee for thoughtful comments that improved this paper, particularly for raising the importance of field geometry effects. This work used the Athena++ MHD code and the DiRAC Complexity system operated by the University of Leicester IT Services, which forms part of the STFC DiRAC HPC Facility (``DiRAC'').

\bibliographystyle{mnras}
\bibliography{references_complete}

\end{document}
